\SourceLeaf{262}
\Article{27}
\indent Ist die gekrümmte Fläche eine Kugel mit dem Radius \(=R\), so gilt
\[
\begin{gathered}
\alpha=\beta=\gamma=-2f^0=\frac{1}{RR};\quad f''=0,\quad g'=0,\\
6h^0-f^0f^0=0\quad\text{oder}\quad k^0=\frac{1}{24R^4}
\end{gathered}
\]
Damit geht Formel [14] über in
\[
A+B+C=\pi+\frac{\sigma}{RR}
\]
Diese Gleichung gilt exakt. Die Formeln 11—13 liefern dagegen
\[
\begin{aligned}
A^*&=A-\frac{\sigma}{3RR}-\frac{\sigma}{180R^4}(2pp-qq+4qq'-q'q')\\
B^*&=B-\frac{\sigma}{3RR}+\frac{\sigma}{180R^4}(pp-2qq+2qq'+q'q')\\
C^*&=C-\frac{\sigma}{3RR}+\frac{\sigma}{180R^4}(pp+qq+2qq'-2q'q')
\end{aligned}
\]
oder mit derselben Genauigkeit
\[
\begin{aligned}
A^*&=A-\frac{\sigma}{3RR}-\frac{\sigma}{180R^4}(bb+cc-2aa)\\
B^*&=B-\frac{\sigma}{3RR}-\frac{\sigma}{180R^4}(aa+cc-2bb)\\
C^*&=C-\frac{\sigma}{3RR}-\frac{\sigma}{180R^4}(aa+bb-2cc)
\end{aligned}
\]
Vernachlässigt man Größen vierter Ordnung, so ergibt sich daraus der bekannte Lehrsatz, den der angesehene \textsc{Legendre} zuerst aufgestellt hat.

\Article{28}
\indent Lässt man die Glieder vierter Ordnung weg, so nehmen unsere allgemeinen Formeln eine sehr einfache Gestalt an
\[
\begin{aligned}
A^*&=A-\tfrac{1}{12}\sigma(2\alpha+\beta+\gamma)\\
B^*&=B-\tfrac{1}{12}\sigma(\alpha+2\beta+\gamma)\\
C^*&=C-\tfrac{1}{12}\sigma(\alpha+\beta+2\gamma)
\end{aligned}
\]
An den Winkeln \(A,B,C\) auf einer nicht kugelförmigen Fläche müssen also unterschiedliche Korrekturen vorgenommen werden, damit die Sinus der korrigierten Winkel zu den gegenüberliegenden Seiten proportional werden. Die Unterschiede zwischen den Korrekturen sind im Allgemeinen von dritter Ordnung. Weicht die Fläche jedoch nur wenig von einer Kugel ab, so gehören diese Unterschiede einer höheren Ordnung an. Selbst bei den größten Dreiecken auf der Erdoberfläche, deren Winkel sich tatsächlich messen lassen, kann der Unterschied stets als
